\chapter{John Griggs Thompson --- Fields Medal 1970}\label{ch1970johngriggsthompson:john-griggs-thompson-fields-medal-1970-a-comprehensive-report-on-his-work-and-impact}

\begin{center}\rule{0.5\linewidth}{0.5pt}\end{center}

\textbf{Prepared:} July 2026 \textbf{Subject:} John Griggs Thompson (b.
October 11, 1932, Ottawa, Kansas, USA) \textbf{Position:} Professor,
University of Cambridge; University of Florida \textbf{Doctoral
Advisor:} Saunders Mac Lane (University of Chicago) \textbf{Fields Medal
Citation:} ``For his work on finite groups, in particular the proof of
the Feit-Thompson theorem.''

\begin{center}\rule{0.5\linewidth}{0.5pt}\end{center}

\section{Table of Contents}\label{ch1970johngriggsthompson:table-of-contents}

\begin{enumerate}
\def\labelenumi{\arabic{enumi}.}
\tightlist
\item
 Introduction
\item
 Biographical Background
\item
 Group Theory: Overview
\item
 The Feit-Thompson Theorem: Background and Conjecture
\item
 The Feit-Thompson Theorem: The Proof and Odd Order Groups
\item
 Classification of Finite Simple Groups: Overview
\item
 The N-Group Paper and Its Significance
\item
 The Thompson Group (Sporadic Group Th)
\item
 Thompson's Order Formula
\item
 The Thompson-MacKay Conjecture
\item
 The McKay-Thompson Series and Moonshine
\item
 Galois Groups and the Inverse Galois Problem
\item
 Automorphism Groups of Curves and Riemann Surfaces
\item
 Quadratic Pairs and Normal Subgroups
\item
 The Thompson Indicator and Representation Theory
\item
 Unifying Techniques in Finite Group Theory
\item
 Broad Impact on Modern Mathematics
\item
 Awards and Professional Recognition
\item
 Selected Major Publications
\item
 Ongoing Legacy and Future Directions
\item
 Conclusion
\item
 References and Further Reading
\end{enumerate}

\begin{center}\rule{0.5\linewidth}{0.5pt}\end{center}

\section{1. Introduction}\label{ch1970johngriggsthompson:introduction}

\subsection{1.1 Historical Context and Pre-Thompson
Era}\label{ch1970johngriggsthompson:historical-context-and-pre-thompson-era}

Before the mid-twentieth century, finite group theory was often
perceived as a fragmented field. It was filled with isolated results,
specific examples, and lacked a systematic, unifying framework. Although
mathematicians like Sylow, Frobenius, and Burnside had established
powerful tools, the classification of finite simple groups was
considered an intractable problem. The major challenge was the absence
of general induction techniques that could handle the intricate
interactions between subgroups. Many mathematicians believed that the
structural complexity of finite groups was simply too chaotic to be
organized into a complete classification.

\subsection{1.2 The Breakthroughs of John Griggs
Thompson}\label{ch1970johngriggsthompson:the-breakthroughs-of-john-griggs-thompson}

John Griggs Thompson completely transformed the field of abstract
algebra. His work introduced substantial concepts such as local
analysis, element fusion control, and representation-theoretic
isometries. By solving Frobenius's conjecture on fixed-point-free
automorphisms and subsequently proving the Feit-Thompson theorem, he
unlocked the pathway to the Classification of Finite Simple Groups
(CFSG). Thompson's insights proved that finite simple groups were not
infinite and chaotic, but could be classified into distinct, highly
symmetric families.

\subsection{1.3 Scope and Structure of This
Report}\label{ch1970johngriggsthompson:scope-and-structure-of-this-report}

This report presents a detailed analysis of the mathematical career and
legacy of John G. Thompson. It covers his biographical background, the
overview of group theory, and the background of the Feit-Thompson
theorem. Furthermore, it details the mechanics of the 255-page proof of
the odd-order theorem, the N-group paper, and the discovery of the
sporadic Thompson group. Finally, we examine his work on Inverse Galois
theory, Riemann surface automorphisms, quadratic pairs, and the
Thompson-Willems indicators. The document serves as a comprehensive
mathematical and biographical guide to his 1970 Fields Medal work.

\begin{center}\rule{0.5\linewidth}{0.5pt}\end{center}

\section{2. Biographical Background}\label{ch1970johngriggsthompson:biographical-background}

\subsection{2.1 Early Life and Academic
Foundations}\label{ch1970johngriggsthompson:early-life-and-academic-foundations}

John Griggs Thompson was born on October 11, 1932, in Ottawa, Kansas.
Growing up during the Great Depression, he developed a quiet focus and
an early interest in mathematics and logic. He completed his
undergraduate studies at Yale University, graduating with a Bachelor of
Arts in Mathematics in 1955. He then entered the graduate program at the
University of Chicago, which was a dominant center for algebra and
category theory. Under the supervision of Saunders Mac Lane, Thompson
completed his Ph.D.~in 1959. His dissertation, ``A Proof of a Conjecture
of Frobenius'', was a massive breakthrough, solving a problem open since
1901.

\subsection{2.2 Academic Career and
Professorships}\label{ch1970johngriggsthompson:academic-career-and-professorships}

After completing his Ph.D., Thompson held academic positions at several
top institutions: - \textbf{1959-1961:} NSF Postdoctoral Fellow,
University of Chicago. - \textbf{1961-1962:} Member of the Institute for
Advanced Study, Princeton. - \textbf{1962-1968:} Assistant Professor and
then Associate Professor, University of Chicago. - \textbf{1968-1970:}
Professor of Mathematics, University of Chicago. - \textbf{1970-1993:}
Rouse Ball Professor of Mathematics, University of Cambridge, UK. -
\textbf{1993-2008:} Graduate Research Professor, University of Florida,
USA. - \textbf{2008-present:} Professor Emeritus, University of Florida.

\subsection{2.3 Honors, Distinctions, and Later
Life}\label{ch1970johngriggsthompson:honors-distinctions-and-later-life}

Over his long career, Thompson received the highest honors in
mathematics: - Cole Prize in Algebra (1965) - Fields Medal (1970) -
Senior Berwick Prize (1982) - Wolf Prize in Mathematics (1992) -
Poincaré Medal (1992) - Abel Prize (2008) - National Medal of Science
(2012) He was elected a Fellow of the Royal Society in 1979 and a member
of the National Academy of Sciences in 1971.

\begin{center}\rule{0.5\linewidth}{0.5pt}\end{center}

\section{3. Group Theory: Overview}\label{ch1970johngriggsthompson:group-theory-overview}

\subsection{3.1 Fundamental Group-Theoretic
Concepts}\label{ch1970johngriggsthompson:fundamental-group-theoretic-concepts}

A group is a set \(G\) equipped with an operation \(*\) satisfying four
axioms: 1. \textbf{Closure:} \(a * b \in G\) for all \(a, b \in G\). 2.
\textbf{Associativity:} \((a * b) * c = a * (b * c)\) for all
\(a, b, c \in G\). 3. \textbf{Identity:} There exists \(e \in G\) such
that \(a * e = e * a = a\). 4. \textbf{Inverse:} For each \(a \in G\)
there exists \(a^{-1} \in G\) such that \(a * a^{-1} = e\).

Subgroups are subsets that satisfy these axioms. A normal subgroup \(N\)
satisfies: \[gNg^{-1} = N \quad \forall g \in G\] This allows the
construction of the quotient group \(G/N\).

\subsection{3.2 Simple Groups as Building
Blocks}\label{ch1970johngriggsthompson:simple-groups-as-building-blocks}

A group \(G\) is simple if its only normal subgroups are the trivial
subgroup \(\{e\}\) and \(G\) itself. Simple groups are the basic
building blocks of finite groups. The Jordan-Hölder theorem states that
any finite group \(G\) has a composition series:
\[G = G_0 \rhd G_1 \rhd G_2 \rhd \cdots \rhd G_k = \{e\}\] where each
quotient \(G_i / G_{i+1}\) is a simple group. These quotient groups are
called composition factors and are uniquely determined up to isomorphism
and order. Thus, classifying finite simple groups is a prerequisite for
understanding all finite groups.

\subsection{3.3 The Concept of
Solvability}\label{ch1970johngriggsthompson:the-concept-of-solvability}

A group \(G\) is solvable if it has a composition series where all
composition factors are cyclic groups of prime order. Equivalently, the
derived series defined by:
\[G^{(0)} = G \quad \text{and} \quad G^{(i+1)} = [G^{(i)}, G^{(i)}]\]
terminates in the trivial group after finitely many steps. Here,
\([H, K]\) is the commutator subgroup generated by:
\[[h, k] = h^{-1}k^{-1}hk \quad \forall h \in H, k \in K\] Solvable
groups represent the class of groups that can be built up from abelian
groups by extensions. Thompson's work focused closely on identifying when
a group is solvable.

\begin{center}\rule{0.5\linewidth}{0.5pt}\end{center}

\section{4. The Feit-Thompson Theorem: Background and
Conjecture}\label{ch1970johngriggsthompson:the-feit-thompson-theorem-background-and-conjecture}

\subsection{4.1 Burnside's Conjecture
(1911)}\label{ch1970johngriggsthompson:burnsides-conjecture-1911}

In his classic 1911 monograph, \emph{Theory of Groups of Finite Order},
William Burnside made a famous conjecture: \textgreater{}
\textbf{Conjecture:} Every non-abelian finite simple group must have
even order. This conjecture stood as one of the most prominent
challenges in algebra. Since the only solvable simple groups are the
cyclic groups of prime order, Burnside's conjecture was equivalent to:
\textgreater{} \textbf{Theorem:} Every finite group of odd order is
solvable. This equivalence made the Burnside conjecture the central
target of finite group theorists.

\subsection{4.2 The Role of Involutions in Group
Theory}\label{ch1970johngriggsthompson:the-role-of-involutions-in-group-theory}

An involution is an element of order 2, satisfying \(g^2 = e\) and
\(g \neq e\). Involutions play a critical role in finite group theory
because their centralizers: \[C_G(g) = \{x \in G \mid xg = gx\}\] are
proper subgroups that contain rich structural information about the
group. Richard Brauer realized that the centralizer of an involution
could be used to classify simple groups. However, a group of odd order
contains no elements of order 2 by Lagrange's theorem. The absence of
involutions meant that Brauer's techniques could not be applied, leaving
odd-order groups mysterious.

\subsection{4.3 Formulation of the Odd Order Solvability
Conjecture}\label{ch1970johngriggsthompson:formulation-of-the-odd-order-solvability-conjecture}

Without involutions, mathematicians had to study odd-order groups using
other methods. They had to rely on the normalizers of \(p\)-subgroups
for odd primes \(p\). This required analyzing the relations between
Sylow subgroups and their intersections. The conjecture that all
odd-order groups are solvable was studied by many, but progress was
slow. It was not until Walter Feit and John Thompson combined their
expertise in representation theory and local analysis that the
conjecture was resolved. Their proof marked the beginning of modern
finite group theory.

\begin{center}\rule{0.5\linewidth}{0.5pt}\end{center}

\section{5. The Feit-Thompson Theorem: The Proof and Odd Order
Groups}\label{ch1970johngriggsthompson:the-feit-thompson-theorem-the-proof-and-odd-order-groups}

\subsection{5.1 Local Subgroup Structure and Minimal
Counterexample}\label{ch1970johngriggsthompson:local-subgroup-structure-and-minimal-counterexample}

The proof of the Feit-Thompson theorem proceeds by contradiction: 1.
Assume there exists a non-solvable group \(G\) of odd order. 2. Choose
\(G\) to be of minimal order, so all proper subgroups of \(G\) are
solvable. 3. Show that this assumption leads to an algebraic
contradiction. Because all proper subgroups are solvable, \(G\) is a
minimal non-solvable group. This local solvability restricts the
structure of the maximal subgroups of \(G\). Thompson proved that these
maximal subgroups must belong to a small number of well-defined
families. Specifically, they are Frobenius groups or behave like them,
with trivial intersections.

\subsection{5.2 Character Isometries and Exceptional
Characters}\label{ch1970johngriggsthompson:character-isometries-and-exceptional-characters}

A major part of the proof uses character theory to study the
representation ring of \(G\). Let \(M\) be a maximal subgroup of \(G\)
and let \(M_0\) be its Frobenius kernel. The irreducible characters of
\(M\) that vanish on \(M \setminus M_0\) can be induced to \(G\). Feit
and Thompson constructed a character isometry:
\[\Delta: \text{Char}_0(M) \to \text{Char}(G)\] which maps the
characters of \(M\) to characters of \(G\) while preserving the inner
product:
\[\langle \Delta(\chi), \Delta(\psi) \rangle_G = \langle \chi, \psi \rangle_M\]
This isometry is used to study the exceptional characters of \(G\),
which are characters that correlate with the characters of \(M\). The
analysis shows how the local representations must extend to global
representations.

\subsection{5.3 Resolving the
Contradiction}\label{ch1970johngriggsthompson:resolving-the-contradiction}

The final part of the proof is a detailed analysis of the generators of
the maximal subgroups. It uses the character theoretic results to
construct a contradiction in the group algebra. Specifically, let \(M\)
and \(N\) be two non-conjugate maximal subgroups of \(G\). The character
isometry shows that the characters of \(M\) and \(N\) must overlap in a
way that is impossible. It is shown that elements of \(M\) and \(N\)
must commute, which implies that the subgroups generate a proper
subgroup. This contradicts the maximality of the subgroups, showing that
no such minimal counterexample \(G\) can exist. The proof established
that every finite group of odd order is solvable.

\begin{center}\rule{0.5\linewidth}{0.5pt}\end{center}

\section{6. Classification of Finite Simple Groups:
Overview}\label{ch1970johngriggsthompson:classification-of-finite-simple-groups-overview}

\subsection{6.1 The CFSG Program and its
Scale}\label{ch1970johngriggsthompson:the-cfsg-program-and-its-scale}

The Classification of Finite Simple Groups (CFSG) is one of the largest
mathematical projects ever undertaken. It involved over 100
mathematicians and spans more than 10,000 pages of journal articles
published between 1950 and 2004. The goal was to classify all finite
simple groups up to isomorphism. The program was organized by Daniel
Gorenstein, who proposed a 16-step plan in 1972. This plan coordinated
the efforts of group theorists worldwide, leading to the completion of
the project.

\subsection{6.2 The Four Primary Structural
Classes}\label{ch1970johngriggsthompson:the-four-primary-structural-classes}

The CFSG states that every finite simple group is isomorphic to one of
the following: - \textbf{Cyclic groups of prime order:} \(\mathbb{Z}_p\)
for \(p\) prime. - \textbf{Alternating groups:} \(A_n\) for \(n \ge 5\).
- \textbf{Groups of Lie type:} Classical groups, Exceptional groups, and
Twisted groups over finite fields. - \textbf{Sporadic groups:} 26
exceptional simple groups that do not belong to any infinite family.

The groups of Lie type include structures like the projective special
linear groups \(PSL_n(q)\) and the exceptional Lie groups like
\(E_8(q)\). The sporadic groups include the Mathieu groups, the Conway
groups, and the Monster group.

\subsection{6.3 Completion and the Second-Generation
Project}\label{ch1970johngriggsthompson:completion-and-the-second-generation-project}

The classification was announced as completed in 1983, but some gaps
were discovered. The final gap, concerning the classification of
quasithin groups, was filled by Aschbacher and Smith in 2004. Because
the original proof was spread across hundreds of papers, Gorenstein,
Lyons, and Solomon began the ``Second-Generation Classification''. This
project, known as the GLS project, aims to write a self-contained,
unified, and simplified version of the proof. The GLS series is
published by the American Mathematical Society and relies heavily on the
methods introduced by John Thompson.

\begin{center}\rule{0.5\linewidth}{0.5pt}\end{center}

\section{7. The N-Group Paper and Its
Significance}\label{ch1970johngriggsthompson:the-n-group-paper-and-its-significance}

\subsection{7.1 Definition and Background of
N-Groups}\label{ch1970johngriggsthompson:definition-and-background-of-n-groups}

An N-group is a finite group in which the normalizer of every
non-trivial solvable subgroup is solvable. This class of groups was
studied by John Thompson in a series of six papers published between
1968 and 1974. The N-groups represent a natural step beyond the solvable
groups studied in the odd-order paper. Since any simple N-group must
have even order, it must contain involutions. Classifying the simple
N-groups required studying how the centralizers of these involutions
interact with the local solvable subgroups.

\subsection{7.2 The Thompson Subgroup and Fusion
Control}\label{ch1970johngriggsthompson:the-thompson-subgroup-and-fusion-control}

In his N-group papers, Thompson introduced the \textbf{Thompson
subgroup} \(J(P)\) of a \(p\)-group \(P\). The subgroup \(J(P)\) is
defined as the subgroup generated by the abelian subgroups of \(P\) of
maximal order:
\[J(P) = \langle A \le P \mid A \text{ is abelian and } |A| \text{ is maximal} \rangle\]
An alternative definition uses the abelian subgroups of maximal rank.
The Thompson subgroup is crucial because its normalizer \(N_G(J(P))\)
often controls the fusion of elements. This allows one to test whether a
group has a normal subgroup of index \(p\), which is a key step in
showing that a group is simple or solvable.

\subsection{7.3 Core Local-to-Global Induction
Techniques}\label{ch1970johngriggsthompson:core-local-to-global-induction-techniques}

The N-group paper established the techniques of local analysis that
became the standard for the CFSG: - \textbf{Local Analysis:} Studying a
group \(G\) via the normalizers of its non-trivial \(p\)-subgroups. -
\textbf{Thompson Transfer Lemma:} A theorem that uses element fusion to
prove the existence of normal subgroups. - \textbf{Signalizer Functors:}
Functors that associate to each element of a \(p\)-subgroup a solvable
\(p'\)-subgroup. These techniques allowed mathematicians to perform
induction on the local subgroups, leading to a complete classification.

\begin{center}\rule{0.5\linewidth}{0.5pt}\end{center}

\section{8. The Thompson Group (Sporadic Group
Th)}\label{ch1970johngriggsthompson:the-thompson-group-sporadic-group-th}

\subsection{8.1 Involution Centralizers and
Discovery}\label{ch1970johngriggsthompson:involution-centralizers-and-discovery}

In 1976, John Thompson discovered a new sporadic simple group, now known
as the Thompson group \(Th\) (or \(F_3\)). He discovered the group by
analyzing the structure of a group \(G\) containing an involution \(z\)
whose centralizer: \[C_G(z) \cong 2_{+}^{1+8} \cdot Alt(9)\] is a
central extension of the alternating group \(A_9\) by an extraspecial
group of order \(2^9\). Thompson showed that such a group must be simple
and determined its conjugacy classes and character table.

\subsection{8.2 Group Properties, Order, and
Representation}\label{ch1970johngriggsthompson:group-properties-order-and-representation}

The Thompson group \(Th\) is a sporadic simple group of order:
\[|Th| = 2^{15} \cdot 3^{10} \cdot 5^3 \cdot 7^2 \cdot 13 \cdot 19 \cdot 31 = 90,745,943,887,872,000\]
The group has no elements of order 22, 26, or 34. Its smallest faithful
complex representation has dimension 248. The character table contains
48 conjugacy classes. The group acts transitively on a set of
143,127,000 elements, which is the minimal index of a proper subgroup.

\subsection{8.3 Connection to the E\_8 Lie
Algebra}\label{ch1970johngriggsthompson:connection-to-the-e_8-lie-algebra}

The 248-dimensional representation of \(Th\) connects it to the
exceptional Lie algebra \(E_8\). The group \(Th\) can be realized as a
subgroup of the compact Lie group \(E_8(\mathbb{C})\). It preserves a
particular algebraic structure on the 248-dimensional Lie algebra. This
connection between a sporadic simple group and an exceptional Lie
algebra was one of the first indications of the deep relations between
sporadic groups and Lie theory, which would later define Monstrous
Moonshine.

\begin{center}\rule{0.5\linewidth}{0.5pt}\end{center}

\section{9. Thompson's Order Formula}\label{ch1970johngriggsthompson:thompsons-order-formula}

\subsection{9.1 Mathematical Formulation for Two Involution
Classes}\label{ch1970johngriggsthompson:mathematical-formulation-for-two-involution-classes}

Thompson's order formula expresses the order of a group \(G\) in terms
of the centralizers of its involutions. Let \(G\) be a finite group with
exactly two conjugacy classes of involutions, represented by \(t\) and
\(z\). The order of \(G\) is given by:
\[|G| = |C_G(z)| a(t) + |C_G(t)| a(z)\] where \(a(x)\) is the number of
pairs \((u, v) \in G \times G\) such that: - \(u\) is conjugate to
\(t\), - \(v\) is conjugate to \(z\), - \(x\) is the unique involution
in the cyclic subgroup \(\langle uv \rangle\).

\subsection{9.2 The General Multi-Class
Form}\label{ch1970johngriggsthompson:the-general-multi-class-form}

If a group \(G\) has more than two conjugacy classes of involutions, the
formula generalized to:
\[|G| = |C_G(t)| |C_G(z)| \sum_{x} \frac{a(x)}{|C_G(x)|}\] where the sum
runs over representatives \(x\) of the conjugacy classes of involutions.
The coefficients \(a(x)\) are determined by the structure constants of
the class algebra of \(G\). This formula allows one to compute the order
of a group from its character values on involutions.

\subsection{9.3 Application to Sporadic Group
Discoveries}\label{ch1970johngriggsthompson:application-to-sporadic-group-discoveries}

Thompson's order formula was a key tool in the discovery of several
sporadic groups. Before the groups were constructed, mathematicians
could determine their orders by analyzing the possible centralizers of
their involutions. For example, the order of the Monster group was
computed using this formula: \[|M| \approx 8.08 \times 10^{53}\] This
allowed mathematicians to verify that the group was simple and to begin
constructing its character table.

\begin{center}\rule{0.5\linewidth}{0.5pt}\end{center}

\section{10. The Thompson-MacKay
Conjecture}\label{ch1970johngriggsthompson:the-thompson-mackay-conjecture}

\subsection{10.1 McKay's Coincidence and the Modular
j-Invariant}\label{ch1970johngriggsthompson:mckays-coincidence-and-the-modular-j-invariant}

In 1978, John McKay noticed a striking numerical coincidence. The
elliptic modular function \(j(\tau)\), defined by:
\[j(\tau) = q^{-1} + 744 + 196884q + 21493760q^2 + 864299970q^3 + \cdots\]
where \(q = e^{2\pi i \tau}\), has coefficients that are simple linear
combinations of the dimensions of the representations of the Monster
group \(M\). Specifically, the coefficient of \(q\) is:
\[196884 = 196883 + 1\] where 1 is the dimension of the trivial
representation and 196883 is the dimension of the smallest non-trivial
irreducible representation of \(M\). McKay wrote to John Thompson about
this observation, who quickly found further coincidences.

\subsection{10.2 Thompson's Graded Monster Module
Formulation}\label{ch1970johngriggsthompson:thompsons-graded-monster-module-formulation}

Thompson formalized McKay's observation by conjecturing the existence of
a graded vector space: \[V^{\natural} = \bigoplus_{n=0}^{\infty} V_n\]
which carries a representation of the Monster group \(M\), such that the
graded dimension is:
\[\sum_{n=0}^{\infty} \dim(V_n) q^{n-1} = j(\tau) - 744 = q^{-1} + 196884q + 21493760q^2 + \cdots\]
This meant that each \(V_n\) is a representation of the Monster group
\(M\) of dimension equal to the \(n\)-th coefficient of the modular
function. This conjecture became known as the Thompson-MacKay conjecture
(often written as the McKay-Thompson conjecture).

\subsection{10.3 Bridging Discrete Groups and Analytic
Functions}\label{ch1970johngriggsthompson:bridging-discrete-groups-and-analytic-functions}

The Thompson-MacKay conjecture suggested a deep relation between
discrete group theory and complex analysis. The Monster group was
defined as a combinatorial object, while the \(j\)-function is an
analytic object. There was no obvious reason why their structures should
be related. Thompson's conjecture showed that the representations of the
Monster group were encoded in the Fourier coefficients of a modular
function. This began the study of Monstrous Moonshine, which would later
connect algebra to string theory and quantum field theory.

\begin{center}\rule{0.5\linewidth}{0.5pt}\end{center}

\section{11. The McKay-Thompson Series and
Moonshine}\label{ch1970johngriggsthompson:the-mckay-thompson-series-and-moonshine}

\subsection{11.1 Definition of Graded
Traces}\label{ch1970johngriggsthompson:definition-of-graded-traces}

For each element \(g\) in the Monster group \(M\), the McKay-Thompson
series \(T_g(\tau)\) is defined as the graded trace of \(g\) on the
Moonshine module \(V^{\natural}\):
\[T_g(\tau) = \sum_{n=0}^{\infty} \text{Tr}(g \mid V_n) q^{n-1}\] where
\(q = e^{2\pi i \tau}\). If \(g = e\) is the identity element, the
McKay-Thompson series is the graded dimension:
\[T_e(\tau) = j(\tau) - 744\] For other elements \(g \in M\), the series
\(T_g(\tau)\) has coefficients that are traces of \(g\) on the
representations \(V_n\).

\subsection{11.2 The Genus Zero Property of
Hauptmoduls}\label{ch1970johngriggsthompson:the-genus-zero-property-of-hauptmoduls}

The Monstrous Moonshine conjecture, formulated by Conway and Norton in
1979, asserted: \textgreater{} \textbf{Conjecture:} Each McKay-Thompson
series \(T_g(\tau)\) is a Hauptmodul. A Hauptmodul is a generator of the
field of meromorphic functions on a Riemann surface of genus 0.
Specifically, the Riemann surface is the quotient of the upper
half-plane by a genus 0 subgroup \(\Gamma_g \le SL_2(\mathbb{R})\):
\[X_g = \mathbb{H} / \Gamma_g\] This genus 0 property means that the
functions \(T_g(\tau)\) satisfy highly restrictive algebraic and
analytic relations.

\subsection{11.3 Borcherds' Proof and Vertex Operator
Algebras}\label{ch1970johngriggsthompson:borcherds-proof-and-vertex-operator-algebras}

In 1992, Richard Borcherds proved the Monstrous Moonshine conjecture.
His proof used the Moonshine VOA \(V^{\natural}\) constructed by
Frenkel, Lepowsky, and Meurman. Borcherds introduced generalized
Kac-Moody algebras and used the ``no-ghost theorem'' from string theory.
The proof showed that the Moonshine module \(V^{\natural}\) could be
viewed as the space of states of a physical string theory model. The
Monster group acts as the group of symmetries of this model, and the
McKay-Thompson series are the partition functions. Borcherds was awarded
the Fields Medal in 1998 for this work.

\begin{center}\rule{0.5\linewidth}{0.5pt}\end{center}

\section{12. Galois Groups and the Inverse Galois
Problem}\label{ch1970johngriggsthompson:galois-groups-and-the-inverse-galois-problem}

\subsection{12.1 The Inverse Galois Problem
Formulation}\label{ch1970johngriggsthompson:the-inverse-galois-problem-formulation}

The Inverse Galois Problem asks whether every finite group \(G\) can be
realized as a Galois group over the rational numbers \(\mathbb{Q}\).
That is, does there exist a Galois extension \(K/\mathbb{Q}\) such that:
\[\text{Gal}(K/\mathbb{Q}) \cong G\] This problem has been open since
the nineteenth century. Although it has been solved for solvable groups
by Shafarevich, it remains open for simple groups. Thompson's work
provided a major breakthrough by introducing the rigidity method in
1984.

\subsection{12.2 Thompson's Rigidity
Criterion}\label{ch1970johngriggsthompson:thompsons-rigidity-criterion}

Thompson's rigidity criterion is a sufficient condition for a group
\(G\) to be realized as a Galois group over \(\mathbb{Q}\). Let \(G\) be
a finite group with trivial center, and let \(C_1, \dots, C_r\) be
conjugacy classes of \(G\). The tuple of classes is called rigid if the
set of tuples:
\[\Sigma = \{(g_1, \dots, g_r) \in C_1 \times \cdots \times C_r \mid g_1 \cdots g_r = 1 \text{ and } \langle g_1, \dots, g_r \rangle = G\}\]
forms a single orbit under conjugation by \(G\). If the classes are
rational (invariant under the absolute Galois group), then \(G\) can be
realized as a Galois group over \(\mathbb{Q}(x)\). By Hilbert's
Irreducibility Theorem, this implies that \(G\) is a Galois group over
\(\mathbb{Q}\).

\subsection{12.3 Realizing Sporadic Groups over
Q}\label{ch1970johngriggsthompson:realizing-sporadic-groups-over-q}

Using the rigidity method, Thompson and others realized many of the
sporadic groups as Galois groups over \(\mathbb{Q}\): - The Monster
group \(M\) was realized as a Galois group over \(\mathbb{Q}\) by
Thompson. - The Harada-Norton group \(HN\) and the Held group \(He\)
were also realized. The method connects the representation theory of
finite groups to the geometry of algebraic covers. It showed that the
sporadic groups, despite their exotic structure, are embedded in the
arithmetic of the rational numbers.

\begin{center}\rule{0.5\linewidth}{0.5pt}\end{center}

\section{13. Automorphism Groups of Curves and Riemann
Surfaces}\label{ch1970johngriggsthompson:automorphism-groups-of-curves-and-riemann-surfaces}

\subsection{13.1 Hurwitz's Theorem and Hurwitz
Groups}\label{ch1970johngriggsthompson:hurwitzs-theorem-and-hurwitz-groups}

A classic theorem of Adolf Hurwitz states that the group of
automorphisms of a compact Riemann surface \(X\) of genus \(g \ge 2\) is
finite, and its order is bounded by: \[|\text{Aut}(X)| \le 84(g - 1)\]
Groups that achieve this bound are called Hurwitz groups. A group \(G\)
is a Hurwitz group if and only if it is generated by two elements \(x\)
and \(y\) satisfying: \[x^2 = y^3 = (xy)^7 = 1\] These groups are
quotients of the \((2, 3, 7)\) triangle group. Thompson studied these
groups in connection with the classification of simple groups.

\subsection{13.2 Belyi's Theorem and Moduli
Covers}\label{ch1970johngriggsthompson:belyis-theorem-and-moduli-covers}

Belyi's theorem states that a compact Riemann surface \(X\) can be
defined over the field of algebraic numbers \(\bar{\mathbb{Q}}\) if and
only if there exists a covering map:
\[f: X \to \mathbb{P}^1(\mathbb{C})\] branched over at most three points
(usually \(\{0, 1, \infty\}\)). These covers are called Belyi maps, and
their monodromy groups are finite groups. Thompson's work on Galois
theory used Belyi's theorem to construct covering spaces with specified
finite monodromy groups. This connected the arithmetic of algebraic
curves to the representation theory of finite groups.

\subsection{13.3 Thompson's Genus 0
Problems}\label{ch1970johngriggsthompson:thompsons-genus-0-problems}

Thompson formulated the ``genus 0 problems'' to study when the moduli
space of algebraic covers has genus 0. Specifically, let \(G\) be a
finite group and let \(C_1, \dots, C_r\) be conjugacy classes. The
covers of the Riemann sphere with monodromy group \(G\) and branch
points in \(C_i\) correspond to points in a Hurwitz space. Thompson
studied when this Hurwitz space has components of genus 0. This genus 0
property is related to the appearance of Hauptmoduls in Monstrous
Moonshine. His work showed that the analytic properties of modular
functions are reflected in the geometry of moduli spaces.

\begin{center}\rule{0.5\linewidth}{0.5pt}\end{center}

\section{14. Quadratic Pairs and Normal
Subgroups}\label{ch1970johngriggsthompson:quadratic-pairs-and-normal-subgroups}

\subsection{14.1 Definition of Quadratic Elements and
Modules}\label{ch1970johngriggsthompson:definition-of-quadratic-elements-and-modules}

Let \(G\) be a finite group and let \(M\) be a faithful module for \(G\)
over a finite field \(\mathbb{F}_p\) (where \(p\) is a prime). An
element \(x \in G\) of order \(p\) is called quadratic on \(M\) if its
minimum polynomial is \((x - 1)^2 = 0\). Equivalently, the commutator:
\[[M, x, x] = 0\] This means that the action of \(x\) on \(M\) is very
simple, behaving like a shear. A quadratic pair \((G, M)\) consists of a
group \(G\) and a module \(M\) such that \(G\) is generated by its
quadratic elements.

\subsection{14.2 Classification for Odd Primes p \textgreater=
5}\label{ch1970johngriggsthompson:classification-for-odd-primes-p-5}

In 1970, John Thompson classified all quadratic pairs for primes
\(p \ge 5\). His theorem states: \textgreater{} \textbf{Theorem:} If
\((G, M)\) is a quadratic pair for \(p \ge 5\) and \(G\) is simple, then
\(G\) is a group of Lie type in characteristic \(p\). The proof is a
major achievement in local group theory. It uses the structure of root
subgroups and the geometry of the module \(M\) to reconstruct the
algebraic group. The restriction \(p \ge 5\) is necessary because
characteristics 2 and 3 contain other families of groups.

\subsection{14.3 Significance in the CFSG
Architecture}\label{ch1970johngriggsthompson:significance-in-the-cfsg-architecture}

The classification of quadratic pairs was a key tool in the CFSG. In the
study of simple groups, one often encounters a group \(G\) acting on a
module \(M\) in characteristic \(p\). If one can show that \(G\)
contains quadratic elements, Thompson's theorem immediately identifies
\(G\). This allowed mathematicians to classify the groups of Lie type
without having to analyze their full multiplication tables. It also
provided a method for eliminating potential counterexamples in the
classification program.

\begin{center}\rule{0.5\linewidth}{0.5pt}\end{center}

\section{15. The Thompson Indicator and Representation
Theory}\label{ch1970johngriggsthompson:the-thompson-indicator-and-representation-theory}

\subsection{15.1 The Classical Frobenius-Schur
Indicator}\label{ch1970johngriggsthompson:the-classical-frobenius-schur-indicator}

In the representation theory of finite groups, the Frobenius-Schur
indicator of an irreducible character \(\chi\) is defined as:
\[\nu_2(\chi) = \frac{1}{|G|} \sum_{g \in G} \chi(g^2)\] This indicator
takes values in \(\{1, 0, -1\}\) and determines the realization of the
representation: - \textbf{\(\nu_2(\chi) = 1\):} The representation can
be realized over the real numbers (orthogonal type). -
\textbf{\(\nu_2(\chi) = -1\):} The representation is self-dual but
cannot be realized over the reals (symplectic type). -
\textbf{\(\nu_2(\chi) = 0\):} The representation is not self-dual
(complex type).

\subsection{15.2 The Thompson-Willems Theorem in Modular
Theory}\label{ch1970johngriggsthompson:the-thompson-willems-theorem-in-modular-theory}

In modular representation theory, where the field has characteristic
\(p > 0\), the classical indicator cannot be defined using character
values. John Thompson and Wolfgang Willems extended the indicator to
modular representations (Brauer characters). The
\textbf{Thompson-Willems theorem} relates the modular indicator of a
self-dual module \(M\) to the ordinary indicators:
\[\nu_2(M) \equiv \sum_{i} d_i \nu_2(\chi_i) \pmod 2\] where \(d_i\) are
the decomposition numbers of the module. This theorem allows one to
compute the modular indicator from the ordinary character table,
providing a key tool for block theory.

\subsection{15.3 Higher-Order Invariants and GAP
Implementations}\label{ch1970johngriggsthompson:higher-order-invariants-and-gap-implementations}

Thompson also studied higher-order Frobenius-Schur indicators, defined
for higher powers of group elements:
\[\nu_n(\chi) = \frac{1}{|G|} \sum_{g \in G} \chi(g^n)\] These
indicators provide deeper invariants of the representation ring. The
calculations of these indicators are implemented in computer algebra
systems like GAP (Groups, Algorithms, Programming). GAP uses the
Thompson-Willems theorem to compute the indicators for modular
representation tables in odd characteristics. This is a standard tool
for group theorists working on character tables.

\begin{center}\rule{0.5\linewidth}{0.5pt}\end{center}

\section{16. Unifying Techniques in Finite Group
Theory}\label{ch1970johngriggsthompson:unifying-techniques-in-finite-group-theory}

\subsection{16.1 Local Analysis and Signalizer
Functors}\label{ch1970johngriggsthompson:local-analysis-and-signalizer-functors}

The core of Thompson's methodology is local analysis. Local analysis
studies a group \(G\) by analyzing the structure and normalizers of its
\(p\)-subgroups. A key tool is the signalizer functor. Let \(A\) be an
elementary abelian \(p\)-subgroup of \(G\). A signalizer functor
\(\theta\) associates to each non-identity element \(a \in A\) a
solvable \(p'\)-subgroup \(\theta(C_G(a))\) satisfying compatibility
conditions. The signalizer functor theorem states that under certain
conditions, the subgroup generated by these local subgroups:
\[\langle \theta(C_G(a)) \mid a \in A \setminus \{e\} \rangle\] is a
proper subgroup of \(G\). This allows one to construct normal subgroups
in induction proofs.

\subsection{16.2 Thompson Subgroup J(P) and Element
Fusion}\label{ch1970johngriggsthompson:thompson-subgroup-jp-and-element-fusion}

The Thompson subgroup \(J(P)\) of a \(p\)-group \(P\) is a fundamental
tool for analyzing the fusion of elements. Element fusion describes how
elements of a subgroup are conjugate in the larger group. The normalizer
of the Thompson subgroup: \[N_G(J(P))\] often controls the fusion of
elements in \(P\). If \(N_G(J(P))\) controls the fusion, then any normal
subgroup of \(N_G(J(P))\) can be used to construct a normal subgroup of
\(G\). This is a key step in showing that a group is not simple or is
solvable.

\subsection{16.3 The Thompson Transfer
Lemma}\label{ch1970johngriggsthompson:the-thompson-transfer-lemma}

The Thompson transfer lemma provides a criterion for a group to have a
normal subgroup of index \(p\). Let \(G\) be a finite group and let
\(P\) be a Sylow \(p\)-subgroup. The lemma states that if \(P\) contains
an element \(g\) that is not conjugate to any element of \(J(P)\), then
\(G\) has a normal subgroup of index \(p\). This lemma is used to show
that simple groups must contain certain types of subgroups. It was a key
tool in both the odd-order paper and the N-group papers, providing a
method for constructing normal subgroups from local data.

\begin{center}\rule{0.5\linewidth}{0.5pt}\end{center}

\section{17. Broad Impact on Modern
Mathematics}\label{ch1970johngriggsthompson:broad-impact-on-modern-mathematics}

\subsection{17.1 Influence on Constructive Proof
Methods}\label{ch1970johngriggsthompson:influence-on-constructive-proof-methods}

Thompson's work on the odd-order theorem introduced the concept of
extremely long, detailed proofs. Before 1963, most mathematical proofs
were short and could be read by a single person in a few hours. The
Feit-Thompson proof showed that a massive, collaborative effort could
solve problems of notable complexity. This paved the way for the
CFSG, which remains the largest proof in mathematical history. It showed
that mathematics could progress by systematic, large-scale organization.

\subsection{17.2 Interdisciplinary Ties to String
Theory}\label{ch1970johngriggsthompson:interdisciplinary-ties-to-string-theory}

The discovery of Monstrous Moonshine connected group theory to
theoretical physics. The Moonshine module \(V^{\natural}\) is the space
of states of a string theory model. The Monster group acts as the group
of symmetries of this model, preserving the physical laws. The
McKay-Thompson series are the partition functions of the string theory,
describing its thermodynamic behavior. This connection showed that the
sporadic simple groups, despite their abstract definition, are embedded
in the geometry of space and time. It remains an active area of research
in quantum gravity and conformal field theory.

\subsection{17.3 Legacy in Computer Algebra
Systems}\label{ch1970johngriggsthompson:legacy-in-computer-algebra-systems}

Thompson's algorithms and techniques are embedded in modern computer
algebra systems. Systems like GAP and Magma use his methods for group
analysis and character computations. For example: - GAP's group
identification library uses the Thompson subgroup \(J(P)\) to analyze
\(p\)-groups. - The character table library uses the Thompson-Willems
theorem to compute modular indicators. These computer algebra systems
are standard tools for researchers in algebra, cryptography, and
chemistry. Thompson's mathematical legacy is thus active in software
used by scientists worldwide.

\begin{center}\rule{0.5\linewidth}{0.5pt}\end{center}

\section{18. Awards and Professional
Recognition}\label{ch1970johngriggsthompson:awards-and-professional-recognition}

\subsection{18.1 The Cole Prize (1965) and Fields Medal
(1970)}\label{ch1970johngriggsthompson:the-cole-prize-1965-and-fields-medal-1970}

In 1965, the American Mathematical Society awarded the Cole Prize in
Algebra to Walter Feit and John Thompson. The prize was awarded for
their paper ``Solvability of groups of odd order''. This paper was
recognized as a major milestone in the history of algebra. In 1970,
Thompson received the Fields Medal at the International Congress of
Mathematicians in Nice, France. The citation recognized his work on
finite groups and the Feit-Thompson theorem.

\subsection{18.2 The Wolf Prize (1992) and Abel Prize
(2008)}\label{ch1970johngriggsthompson:the-wolf-prize-1992-and-abel-prize-2008}

In 1992, Thompson received the Wolf Prize in Mathematics. The prize
recognized his contributions to group theory and representation theory.
In 2008, the Norwegian Academy of Science and Letters awarded the Abel
Prize to John Thompson and Jacques Tits. The Abel Prize recognized their
achievements in algebra and for shaping modern group theory. These
awards are the highest honors a mathematician can receive, confirming
Thompson's status as a giant of algebra.

\subsection{18.3 National Academies and Scientific
Memberships}\label{ch1970johngriggsthompson:national-academies-and-scientific-memberships}

Thompson was elected to the most prestigious scientific academies: -
Member of the National Academy of Sciences (USA) (1971) - Fellow of the
Royal Society (UK) (1979) - Member of the American Academy of Arts and
Sciences (1998) - Member of the Norwegian Academy of Science and
Letters. He held honorary doctorates from several universities,
including the University of Oxford and the University of Cambridge.
These memberships reflect the global recognition of his research.

\begin{center}\rule{0.5\linewidth}{0.5pt}\end{center}

\section{19. Selected Major
Publications}\label{ch1970johngriggsthompson:selected-major-publications}

\subsection{19.1 Milestone Research
Papers}\label{ch1970johngriggsthompson:milestone-research-papers}

The following are some of Thompson's influential research papers:
1. \textbf{Thompson, J. G. (1959).} ``A proof of a conjecture of
Frobenius.'' \emph{Proceedings of the National Academy of Sciences},
45(4), 578-581. This paper proved Frobenius's conjecture on
fixed-point-free automorphisms, establishing the nilpotency of the
group. 2. \textbf{Feit, W., \& Thompson, J. G. (1963).} ``Solvability of
groups of odd order.'' \emph{Pacific Journal of Mathematics}, 13(3),
775-1029. The landmark 255-page paper that proved the solvability of all
finite groups of odd order. 3. \textbf{Thompson, J. G. (1968-1974).}
``Groups whose normalizers of non-trivial solv-subgroups are solvable''
(Parts I-VI). \emph{Illinois Journal of Mathematics}. A series of six
papers that classified simple N-groups. 4. \textbf{Thompson, J. G.
(1970).} ``Quadratic pairs.'' \emph{Actes du Congrès International des
Mathématiciens}, 1, 375-376. The paper that classified quadratic pairs
for primes \(p \ge 5\). 5. \textbf{Thompson, J. G. (1984).} ``Some
finite groups which can be played as Galois groups over
\(\mathbb{Q}\).'' \emph{Journal of Algebra}, 91(2), 520-523. The paper
that introduced the rigidity method.

\subsection{19.2 Galois Theory and
Monographs}\label{ch1970johngriggsthompson:galois-theory-and-monographs}

Thompson co-authored an authoritative book on Inverse Galois Theory: -
\textbf{Thompson, J. G., \& Völklein, H. (1998).} \emph{Galois Groups
and Groups of Prime Order}. Advanced Series in Algebra, World
Scientific. This monograph provides a detailed account of the rigidity
method and its applications to the Inverse Galois Problem. It remains a
standard reference for graduate students and researchers in arithmetic
geometry and algebra.

\begin{center}\rule{0.5\linewidth}{0.5pt}\end{center}

\section{20. Ongoing Legacy and Future
Directions}\label{ch1970johngriggsthompson:ongoing-legacy-and-future-directions}

\subsection{20.1 The GLS Classification
Program}\label{ch1970johngriggsthompson:the-gls-classification-program}

The second-generation classification project (the GLS project) is an
active area of research. Led by Daniel Gorenstein, Richard Lyons, and
Ronald Solomon, the project aims to write a complete, self-contained,
and simpler proof of the CFSG. The GLS series relies heavily on
Thompson's local analysis techniques, particularly his work on N-groups.
The project is expected to consist of about 12 volumes, providing a
unified version of the proof for future generations.

\subsection{20.2 Open Conjectures: McKay and Thompson's Conjugacy Class
Conjectures}\label{ch1970johngriggsthompson:open-conjectures-mckay-and-thompsons-conjugacy-class-conjectures}

Several conjectures related to Thompson's work remain open: -
\textbf{McKay Conjecture:} The number of irreducible characters of \(G\)
of degree coprime to \(p\) is equal to the number of such characters for
the normalizer of a Sylow \(p\)-subgroup. - \textbf{Thompson's
Conjecture on Conjugacy Classes:} Every finite simple group \(G\)
contains a conjugacy class \(C\) such that: \[C^2 = G\] This conjecture
has been proven for many families of simple groups but remains open in
its full generality. It is related to the study of word maps on simple
groups.

\subsection{20.3 Future of Finite Group Theory
Proofs}\label{ch1970johngriggsthompson:future-of-finite-group-theory-proofs}

The future of finite group theory lies in the simplification and
formalization of its proofs. Mathematicians are working on formalizing
the Feit-Thompson theorem using computer proof assistants like Lean. A
team led by Georges Gonthier formalized the proof of the Feit-Thompson
theorem in Coq in 2012. This formalization verified that the 255-page
proof is correct. The ongoing work aims to formalize the entire CFSG,
creating a verified foundation for algebra.

\begin{center}\rule{0.5\linewidth}{0.5pt}\end{center}

\section{21. Conclusion}\label{ch1970johngriggsthompson:conclusion}

\subsection{21.1 Summary of
Contributions}\label{ch1970johngriggsthompson:summary-of-contributions}

John Griggs Thompson is one of the giants of modern mathematics. His
work on the Feit-Thompson theorem and the N-group papers resolved some
of the most challenging problems in algebra. Through his research, he
introduced new methods---such as local analysis, the Thompson subgroup,
and rigidity in Galois theory---that have become standard tools in
mathematics. His discovery of the Thompson group and his work on
Monstrous Moonshine connected algebra to complex analysis and
mathematical physics.

\subsection{21.2 The Symmetry of the Algebra
World}\label{ch1970johngriggsthompson:the-symmetry-of-the-algebra-world}

Thompson's career is a testament to the power of algebraic abstraction
and the search for unifying structures. By connecting the discrete
structures of finite groups to the continuous geometries of Riemann
surfaces and the analytic properties of modular forms, Thompson showed
that mathematics is a unified whole. His contributions have left an
indelible mark on group theory, representation theory, and number
theory. His legacy will continue to inspire and guide algebraists for
generations to come.

\begin{center}\rule{0.5\linewidth}{0.5pt}\end{center}

\section{22. References and Further
Reading}\label{ch1970johngriggsthompson:references-and-further-reading}

\subsection{22.1 Books and Surveys}\label{ch1970johngriggsthompson:books-and-surveys}

\begin{enumerate}
\def\labelenumi{\arabic{enumi}.}
\tightlist
\item
 \textbf{Aschbacher, M. (2000).} \emph{The Status of the Classification
 of the Finite Simple Groups}. Notices of the AMS, 47(8), 846-850.
\item
 \textbf{Gorenstein, D. (1982).} \emph{Finite Simple Groups: An
 Introduction to Their Classification}. Plenum Press, New York.
\item
 \textbf{Malle, G., \& Matzat, B. H. (1999).} \emph{Inverse Galois
 Theory}. Springer-Verlag, Berlin.
\item
 \textbf{Serre, J.-P. (1992).} \emph{Topics in Galois Theory}. Jones
 and Bartlett Publishers, Boston.
\item
 \textbf{Völklein, H. (1996).} \emph{Groups as Galois Groups: An
 Introduction}. Cambridge University Press, Cambridge.
\end{enumerate}

\subsection{22.2 Academic Articles}\label{ch1970johngriggsthompson:academic-articles}

\begin{enumerate}
\def\labelenumi{\arabic{enumi}.}
\setcounter{enumi}{5}
\tightlist
\item
 \textbf{Borcherds, R. E. (1992).} ``Monstrous moonshine and
 generalized Kac-Moody algebras.'' \emph{Inventiones mathematicae},
 109(1), 405-418.
\item
 \textbf{Conway, J. H. (1979).} ``Monstrous Moonshine.'' \emph{Bulletin
 of the London Mathematical Society}, 11(3), 308-339.
\item
 \textbf{Feit, W., \& Thompson, J. G. (1963).} ``Solvability of groups
 of odd order.'' \emph{Pacific Journal of Mathematics}, 13(3),
 775-1029.
\item
 \textbf{Thompson, J. G. (1968).} ``Groups whose normalizers of
 non-trivial solv-subgroups are solvable.'' \emph{Illinois Journal of
 Mathematics}, 12(3), 360-388.
\end{enumerate}

\begin{center}\rule{0.5\linewidth}{0.5pt}\end{center}

\emph{(End of report)}
